\medskip

L'objectif de cet exercice est de conjecturer en partie A puis de démontrer en partie B le comportement d'une suite.

Les deux parties peuvent cependant être traitées de manière indépendante.

On considère la suite $\left(u_n\right)$ définie par $u_0 = 3$ et pour tout $n \in \N$ :

\[u_{n+1} = \dfrac{4}{5 - u_n}.\]

\medskip

\textbf{Partie A}

\medskip

\begin{enumerate}
\item Recopier et compléter la fonction Python suivante \texttt{suite(n)} qui prend comme
paramètre le rang $n$ et renvoie la valeur du terme $u_n$.

\begin{center}
{\ttfamily
\begin{tabular}{|l|}\hline
\pmb{def} suite(n):\\
\qquad u = \ldots\\
\qquad \pmb{for} i \pmb{in} range(n) : \quad \\
\qquad \qquad \ldots\\ 
\qquad \pmb{return} u \\ \hline
\end{tabular}
}
\end{center}

\item L'exécution de \texttt{suite(2)} renvoie 1.3333333333333333. 

Effectuer un calcul pour vérifier et expliquer cet affichage.
\item À l'aide des affichages ci-dessous, émettre une conjecture sur le sens de variation et une conjecture sur la convergence de la suite $\left(u_n\right)$.

\begin{center}
{\ttfamily
\begin{tabular}{|l|}\hline
>> suite(2) \\
1.3333333333333333\\
>>  suite(5)\\
1.0058479532163742\\
>>  suite(10)\\
1.0000057220349845\\
>>  suite(20)\\
1.000000000005457\\ \hline
\end{tabular}
}
\end{center}

\end{enumerate}

\medskip

\textbf{Partie B}

\medskip

On considère la fonction $f$ définie et dérivable sur l'intervalle $] -\infty~;~5[$ par:

\[f(x) = \dfrac{4}{5 - x}.\]

Ainsi, la suite $\left(u_n\right)$ est définie par $u_0 = 3$ et pour tout $n \in \N,\:\: u_{n+1}  = f\left(u_n\right)$.

\medskip

\begin{enumerate}
\item Montrer que la fonction $f$ est croissante sur l'intervalle $]-\infty~;~5[$.
\item Démontrer par récurrence que pour tout entier naturel $n$ on a : 
\[1 \leqslant u_{n+1} \leqslant u_n \leqslant 4.\]
\item 
	\begin{enumerate}
		\item Soit $x$ un réel de l'intervalle $]-\infty~;~5[$.

Prouver l'équivalence suivante: 

\[f(x) = x \iff x^2 - 5x + 4 = 0.\]
		
		\item Résoudre $f(x) = x$ dans l'intervalle $]-\infty~;~5[$.
	\end{enumerate}	
\item Démontrer que la suite $\left(u_n\right)$ est convergente.

Déterminer sa limite.
\item Le comportement de la suite serait-il identique en choisissant comme terme initial $u_0 = 4$ au lieu de $u_0 = 3$ ?
\end{enumerate}
